% ---
% id: meat-core/prop:drop-bound
% kind: prop
% title: "The dropped program bounds the value from below"
% summary: >-
%   Setting the deficits of a dropped set to zero can only lower the
%   objective, so the dropped value lies between 0 and h_low, hence below the
%   single-round rate, and never grows as the set grows. Dropping only
%   key-inert pairs changes nothing.
% prominence: standard
% status: accepted
% uses: [def:dropped, def:joint-sectors, lem:joint-sectors, def:feasible, foundations/lem:kernel, cor:reduction-constrained, foundations/def:kl, def:admissible-declaration]
% mentions: []
% assumes: []
% refs: []
% provenance: newly-derived
% lean:
% ---

Deleting summands that are all non-positive can only raise each $K$, and
the objective falls as $K$ rises. The loss is in the rate and needs no
hypothesis on the key map or the key holder, and it is a bound only because
the dropped sectors stay in the program.

\begin{proposition}\label{prop:drop-bound}
Let $\Decl$ satisfy Definition~\ref{def:admissible-declaration} in the general format, in
either source mode and with an arbitrary list $\CList$; let
$\alpha\in(1,\infty)$ and let $\mathcal{D}\subseteq\mathcal{D}'$ be dropped
sets (Definition~\ref{def:dropped}). For every $\omega$ block diagonal along
the joint sectors, every $\hat c\in\Chkey$ and every distribution $\mathbf q$
on $\Chatset$:
\begin{enumerate}[label=(\roman*),leftmargin=2.2em]
\item $0\le K_{\hat c}(\omega)\le K^{\mathcal{D}}_{\hat c}(\omega)
\le K^{\mathcal{D}'}_{\hat c}(\omega)\le P_{\hat c}(\omega)$;
\item $O^{\mathcal{D}'}_{\alpha}(\mathbf q,\omega)\le
O^{\mathcal{D}}_{\alpha}(\mathbf q,\omega)\le O_{\alpha}(\mathbf q,\omega)$
in the extended-real sense, and $O^{\mathcal{D}'}_{\alpha}(\mathbf
q,\omega)\ge0$ whenever $\tr\omega=1$;
\item consequently
$0\le\underline{h}^{\mathcal{D}'}_{\alpha}\le
\underline{h}^{\mathcal{D}}_{\alpha}\le\hlow\le\hqkd(\CList)$;
\item if every pair of $\mathcal{D}$ is key-inert
(Definition~\ref{def:joint-sectors}) then
$K^{\mathcal{D}}_{\hat c}=K_{\hat c}$ and
$O^{\mathcal{D}}_{\alpha}=O_{\alpha}$ identically.
\end{enumerate}
\end{proposition}

\begin{proof}
(i) By \eqref{eq:K-joint} and \eqref{eq:KD}, the three quantities
$K_{\hat c}$, $K^{\mathcal{D}}_{\hat c}$, $K^{\mathcal{D}'}_{\hat c}$ are
$P_{\hat c}(\omega)$ plus the sum of $f_{\hat c,(t,s)}(\omega^{(t,s)})$ over
all joint sectors, over those with $(\hat c,(t,s))\notin\mathcal{D}$, and
over those with $(\hat c,(t,s))\notin\mathcal{D}'$ respectively: three
nested index sets, the last the smallest. Every summand is $\le0$
(Lemma~\ref{lem:joint-sectors}\,(ii)), so deleting summands does not
decrease the sum, which gives the two middle inequalities, and the last
holds because the remaining sum is $\le0$. $K_{\hat c}(\omega)\ge0$: by
Lemma~\ref{lem:joint-sectors}\,(ii),
$P_{\hat c,(t,s)}+f_{\hat c,(t,s)}\ge0$ on every joint sector, and
$K_{\hat c}(\omega)$ is the sum of these over the joint sectors.

(ii) In \eqref{eq:objective} the quantity $K$ enters only through the terms
$\frac{1}{\alpha-1}d\big(q(\hat c),K\big)$, $\hat c\in\Chkey$, with $d$ the
kernel of \xref{foundations/lem:kernel} --- whose conventions at vanishing arguments
are those of \xref{foundations/def:kl} --- and the prefactor is positive. $d$
is non-increasing in its second argument, so (i) gives the two inequalities
termwise, in the extended-real sense. For the lower bound let $\tr\omega=1$,
so that $\mathbf P\coloneqq\Phimap[\omega]$ is a probability distribution. If
some $q(\hat c)>0$ has $P_{\hat c}=0$, or
$K^{\mathcal{D}'}_{\hat c}=0$ with $\hat c\in\Chkey$, the objective is
$+\infty$. Otherwise, by $K^{\mathcal{D}'}_{\hat c}\le P_{\hat c}$ and
monotonicity of $d$,
$d(q(\hat c),K^{\mathcal{D}'}_{\hat c})\ge d(q(\hat c),P_{\hat c})$, so
\begin{equation*}
  O^{\mathcal{D}'}_{\alpha}(\mathbf q,\omega)\;\ge\;
  \frac{\alpha}{\alpha-1}\sum_{\hat c\in\Chatset}
  q(\hat c)\log_2\frac{q(\hat c)}{P_{\hat c}}
  =\frac{\alpha}{\alpha-1}\,\sfD(\mathbf q\,\|\,\mathbf P)\;\ge\;0 ,
\end{equation*}
the terms $\hat c\notin\Chkey$ carrying the prefactor
$\frac{\alpha}{\alpha-1}$ already and the terms $\hat c\in\Chkey$ the sum
$1+\frac{1}{\alpha-1}$ of their two prefactors; the last step is
non-negativity of the divergence between probability distributions,
Gibbs' inequality of \xref{foundations/def:kl}.

(iii) Take infima over $\acSet\times\Feas$ in (ii); every $\omega\in\Feas$
has $\tr\omega=1$. The last inequality is
Corollary~\ref{cor:reduction-constrained}.

(iv) For a key-inert pair $f_{\hat c,(t,s)}\equiv0$
(Lemma~\ref{lem:joint-sectors}\,(iii)), so the summands deleted in
\eqref{eq:KD} are identically zero.
\end{proof}
